% !TEX root = main.tex
% ---------------------------------------------------------------------------
% Method section — assumes the v3.3 synthesis recipe and a trained generalist.
% Written to be venue-agnostic: Sec. 3.1-3.8 are shared by all three framings in
% outlines.tex; only the emphasis (which subsection is "the contribution") changes.
% ---------------------------------------------------------------------------
\section{Method}
\label{sec:method}

We describe a pipeline that produces a gentle grasping policy for deformable, fragile
objects without any contact sensing at deployment. The organising idea is that the
quantity which defines gentleness---internal stress---is unobservable to the robot but
\emph{exact} inside a deformable simulator. We therefore spend it where it exists: an
autonomous demonstrator selects grasps by minimising predicted internal stress
(\S\ref{sec:method:synth}), executes them under a randomised scene distribution
(\S\ref{sec:method:demos}), and a point-cloud policy is behaviour-cloned from the result
(\S\ref{sec:method:policy}) and deployed on hardware from depth and proprioception alone
(\S\ref{sec:method:deploy}).

\subsection{Problem setting}
\label{sec:method:problem}

A 7-DoF arm with a parallel-jaw gripper must lift a single fragile object from a table and
hold it, without damaging it. The object's pose, size, shape and material stiffness vary
between episodes and are not directly observable. Perception is a single fixed external
depth camera; there is no tactile or force sensing in the loop.

Demonstrations are generated in a material-point-method (MPM) simulator that represents the
object as a deformable continuum, so that internal stress is available at every particle and
the contact between the gripper and the object is resolved rather than scripted. The same
task, workspace and camera geometry are instantiated on hardware, and the boundary between
the two is a single observation struct: everything above it (point-cloud processing, action
encoding) is shared code, everything below it is backend-specific.

At each control step $t$ the policy observes
\begin{equation}
  o_t \;=\; \bigl(\, \mathbf{p}_t,\; \mathbf{q}_t,\; w_t,\; \mathcal{P}_t \,\bigr),
  \label{eq:obs}
\end{equation}
the end-effector position $\mathbf{p}_t\!\in\!\mathbb{R}^3$, its orientation as a unit
quaternion $\mathbf{q}_t\!\in\!S^3$, the measured gripper width $w_t\!\in\!\mathbb{R}$, and a
point cloud $\mathcal{P}_t\!\in\!\mathbb{R}^{N_p\times 3}$ with $N_p\!=\!1024$ points in the
robot base frame. The action is an \emph{absolute} target pose and gripper width,
\begin{equation}
  a_t \;=\; \bigl(\, \mathbf{p}^{\mathrm{cmd}}_t,\; \boldsymbol{\theta}^{\mathrm{cmd}}_t,\;
                    w^{\mathrm{cmd}}_t \,\bigr) \;\in\; \mathbb{R}^{7},
  \label{eq:act}
\end{equation}
with orientation encoded as Euler angles (\S\ref{sec:method:spaces}). An episode succeeds if
the object's centre enters a height band and remains there for $T_{\text{hold}}$ consecutive
steps; it is \emph{gentle} to the degree that the peak internal (von~Mises) stress stays
below the material's bruising threshold. Success is measurable on hardware; gentleness, in
this work, is measured in simulation.

One design decision follows from this asymmetry and is used throughout: gentleness is
enforced at \emph{data-generation} time by the demonstrator rather than by a reward term or a
feedback controller, so the deployed policy never needs a signal it cannot observe.

\subsection{Stress-aware grasp synthesis}
\label{sec:method:synth}

The demonstrator plans, for a given object instance, a $7$-DoF grasp
$x = (\mathbf{t}, \mathbf{r}, \omega)$: the tool-centre-point position $\mathbf{t}$, its
orientation $\mathbf{r}$ (Euler), and the \emph{commanded} gripper width $\omega$. Planning
happens on the mesh, outside the physics simulator.

\paragraph{Width-controlled contact model.}
A real parallel jaw is position-controlled: one commands a width and the grip force is an
outcome. We model this directly. The object is tetrahedralised and a linear-elastic
stiffness matrix $K$ assembled. Because the object is free-floating, $K$ is rank-deficient
by the six rigid-body modes $R$, so all solves use an inertia-relief (bordered) system
factored once per object,
\begin{equation}
  \begin{bmatrix} K & R \\ R^{\!\top} & 0 \end{bmatrix}
  \begin{bmatrix} u \\ \alpha \end{bmatrix}
  =\begin{bmatrix} b \\ 0 \end{bmatrix}.
  \label{eq:inertia-relief}
\end{equation}
Closing to width $\omega$ indents the object by $\delta_{l},\delta_{r}$ per jaw. We prescribe
\emph{normal-only} displacement at the contact nodes---the tangential directions stay
free---with a rounded (parabolic) pad profile $d_i=\delta\,\bigl(1-(r_i/r_{\text{pad}})^2\bigr)$.
Both details matter: fully-bonded contact against a flat plane clamps the surface and
produces a mesh-dependent flat-punch edge singularity instead of bulk compression, whereas
the normal-only rounded pad lets the material bulge and yields a physical compression zone.

Under position control the displacement field is independent of Young's modulus, so stress
and grip force are both \emph{linear} in it. We solve once at $E\!=\!1$ to obtain
$(\sigma_1, F_1)$ and recover any material by scaling,
\begin{equation}
  \sigma(E) \;=\; E\,\sigma_1, \qquad F(E) \;=\; E\,F_1 ,
  \label{eq:E-linear}
\end{equation}
which makes evaluating a candidate grasp across a distribution of stiffnesses essentially
free---one FEM solve per pose, scalar arithmetic per material sample. A grasp is
\emph{holdable} when friction at the two pads carries the object's weight with an
acceleration margin $a$,
\begin{equation}
  2\,\mu\,F(E) \;\ge\; m\,g\,\bigl(1 + a/g\bigr).
  \label{eq:hold}
\end{equation}

\paragraph{Grasp score.}
Among holdable grasps we prefer those that deform the object least, press flush, and do not
concentrate load on a small patch. Writing $\sigma_{\text{top10}}$ for the mean of the
highest decile of element stresses (with elements adjacent to the contacts masked, since a
point load produces a mesh-dependent spike there) and $\sigma_{p98}$ for the unmasked $98$th
percentile, the objective to maximise is
\begin{equation}
  \begin{aligned}
  S(x) \;=\; &-\,\sigma_{\text{top10}}(x)
              \;-\; w_{\text{al}}\bigl(1-\mathrm{al}(x)\bigr)
              \;-\; w_{\text{pk}}\,\sigma_{p98}(x) \\[2pt]
             &-\; w_{\text{pr}}\,\frac{F(x)}{A_{\min}(x)}
              \;-\; w_{\text{az}}\,\bigl[\theta_{\text{az}}(x)-\bar\theta\bigr]_{+},
  \end{aligned}
  \label{eq:score}
\end{equation}
subject to holdability~\eqref{eq:hold} and a hard floor on the contact area of the
\emph{worst} pad, $A_{\min}(x) \ge A_{\text{floor}}$; infeasible candidates are returned on a
shaped penalty ladder so the search retains a gradient toward feasibility. The terms encode
distinct failure modes we observed. The alignment term $\mathrm{al}(x)$, the mean absolute
cosine between the closing axis and the local surface normals, rejects grazing and edge
contacts. The peak and pressure terms penalise \emph{local} damage that the masked bulk
statistic hides. The area floor is what eliminates stem-type grasps: on our object these
score \emph{better} on bulk stress than a correct enveloping grasp (squeezing near a stiff,
slender part deforms the body less), so a search that optimises bulk stress alone converges
toward them; they are separated instead by worst-pad contact area, which differs by roughly
a factor of six between the two grasp types. Finally $\theta_{\text{az}}$ is the azimuth of
the closing axis relative to the camera ray, bounded by $\bar\theta$ so the fingers do not
occlude the object in the very observation the policy must act on. A shaped penalty rather
than a hard rejection is required here: occlusion-reducing candidates otherwise sit on a flat
infeasibility floor where a soft weight has no gradient.

\paragraph{Search.}
$S$ is optimised with CMA-ES from multiple seeds---canonical closing axes plus a yaw fan
centred on the camera-perpendicular direction---since a single start reliably misses the
gentle basin. The FEM factorisation is built once per object geometry and reused across all
candidates and all envs sharing that geometry, so per-candidate cost is one Schur-complement
back-substitution.

\subsection{Demonstration generation}
\label{sec:method:demos}

Planned grasps are executed in a deformable (MPM) simulator to produce
$(\text{observation},\text{action})$ trajectories in the same format a real teleoperation
session would produce.

\paragraph{Execution.}
Each parallel environment runs an independent phase machine---approach, settle, close,
firm, lift, hold---so that environments may progress at different rates and, on failure,
rewind. Commands are still issued as one batched step; only the per-environment row varies.
At the close$\rightarrow$firm boundary the measured contact signal is read once: a grasp that
came out weak receives an additional closure, which recovers marginal grasps without
introducing over-squeeze elsewhere. Slip during the lift triggers a bounded regrasp, and the
failed attempt is retained in the demonstration by design, so recovery behaviour is
represented in the data.

\paragraph{Approach trajectory.}
Rather than interpolating linearly from home to the grasp pose, the approach reproduces the
structure we measured in human teleoperation: the operator aligns laterally while still well
above the object, then descends, with orientation adjusting throughout. We implement this as
\emph{per-axis progress} over a single phase, with no waypoint and no stop:
\begin{equation}
  s_{xy}(\alpha) \;=\; \mathrm{smoothstep}\!\Bigl(\min\bigl(\tfrac{\alpha}{f},\,1\bigr)\Bigr),
  \qquad
  s_{z}(\alpha) \;=\; \alpha ,
  \qquad \alpha \in (0,1],
  \label{eq:approach}
\end{equation}
where $f\!\sim\!\mathcal{U}[f_{\text{lo}},f_{\text{hi}}]$ is drawn per environment and sets
how early the lateral motion completes; orientation is slerped over the full $\alpha$. The
lateral velocity reaches zero only at its own arrival, while $z$ is still moving, so the path
is continuous and never dwells (\S\ref{sec:method:supervision} explains why that matters).
The phase \emph{duration} is set from the profile's arc length $L(f)$ and a reference speed,
\begin{equation}
  T_{\text{app}} \;=\; \mathrm{clip}\bigl(\,L(f)/v_{\text{ref}},\; T_{\min},\, T_{\max}\bigr),
  \label{eq:duration}
\end{equation}
so that per-step speed is approximately constant across episodes. A fixed duration would
instead make speed proportional to the spawn distance; measured on our demonstrations this
produces a speed--distance correlation far outside what human teleoperation exhibits, and it
makes the commanded action magnitude depend on where the episode started rather than on the
current state.

\paragraph{Scene randomisation.}
Each scene draws object pose in the physical workspace rectangle, full yaw and bounded
pitch/roll with a probability of a large flip (so the object may settle inverted), uniform
scale, procedural mesh deformation (bend, twist, taper, anisotropic axis scale), material
($E$, $\nu$, $\rho$, coupling friction), a base mesh from a per-category pool, and a small
jitter on the arm's home pose. Geometry-changing draws require a scene rebuild and are
applied at a configurable cadence; pose-level draws are per reset.

\paragraph{Post-processing.}
Two dataset-level steps follow collection. First, runs of identical held commands longer
than a threshold are collapsed to a fixed number of frames; this is applied only to absolute
actions, where a held command is a literal repeated target, and the retained count at the end
of the episode is what supervises \emph{stopping} once the lift height is reached. Second, a
geometric filter removes episodes whose final configuration indicates the object dangling
from the fingertips rather than being enveloped---diagnosed from the tool-centre-point
position relative to the object centre at hold, with the width and lateral-offset cues
applied conditionally so that genuinely small objects are not misclassified.

\subsection{Observation and action spaces}
\label{sec:method:spaces}

\paragraph{Perception.}
Depth is back-projected with the camera intrinsics and extrinsics into the robot base frame,
cropped to the workspace, denoised by a voxel-occupancy outlier filter (which is a no-op on
clean simulated depth and removes flying pixels on the real sensor), and reduced to $N_p$
points by farthest-point sampling. Because the arm occupies most of the workspace, an
\emph{arm-focus} reweighting samples points that are neither low nor near the end-effector at
a reduced rate, so the fixed budget concentrates on the object and the fingers. This
reweighting is applied identically in both domains---it is scene geometry, not sensor
noise---whereas the outlier filter is deliberately domain-specific. Quaternions are
sign-canonicalised (largest-magnitude component positive) in the shared pipeline, which is a
no-op in simulation and removes a systematic sign difference reported by the real arm.

\paragraph{Actions.}
The raw policy output is clipped to $[-1,1]^7$ and mapped affinely onto workspace position
bounds, Euler ranges and the physical gripper range. The Euler encoding carries a constant
frame offset chosen so that the nominal downward-pointing tool orientation lies away from the
$\pm\pi$ wrap, removing a seam that otherwise appears in the middle of the action
distribution. The decoded absolute target is then rate-limited against the \emph{previous
commanded} target,
\begin{equation}
  \Delta \mathbf{p} \leftarrow \mathrm{clip}(\Delta \mathbf{p}, \pm\rho_{p}), \quad
  \Delta \mathbf{r} \leftarrow \mathrm{clip}(\log(R_{\text{cmd}}R_{\text{prev}}^{\!\top}), \pm\rho_{r}), \quad
  \Delta w \leftarrow \mathrm{clip}(\Delta w, \pm\rho_{w}),
  \label{eq:rate}
\end{equation}
with the rotation clipped per axis in the world rotation-vector representation. The same
bound is enforced in the demonstration generator, so the dataset is bounded by construction,
and at both backends, so an out-of-distribution policy output executes as a bounded walk
rather than a full-speed servo jump. Absolute (rather than delta) actions are used because
they re-anchor to a pose every step and therefore do not accumulate execution error across a
chunk.

\subsection{Action supervision: lead and dwell}
\label{sec:method:supervision}

Behaviour cloning of absolute-pose actions has two failure modes that produce a policy with
excellent training loss and no useful behaviour. We state them as dataset conditions because
both are properties of the demonstrations, not of the architecture.

The first concerns \emph{what} is recorded. If the action label is the pose the robot
\emph{achieved} at the next step, then in a closed loop the policy can satisfy its objective
by commanding where it already is: since the controller lags its target, achieved poses trail
commands, and the learned map has a fixed point at $a_t \approx x_t$. Supervision must use
targets that lead the state---in practice the recorded \emph{commanded} targets, which carry
the controller's tracking lag inherently, or a $K$-step lookahead if only achieved poses
survive. We express the requirement as a dataset predicate on the lead distribution,
\begin{equation}
  \mathrm{quantile}_{75}\bigl(\lVert a_t - x_t \rVert\bigr) \;\ge\; \varepsilon_{\text{lead}} .
  \label{eq:lead}
\end{equation}

The second concerns \emph{how the trajectory moves}. Even with commanded supervision, a
demonstration whose velocity passes through zero mid-trajectory---the natural consequence of
minimum-jerk time scaling, and of holding a pose---presents the policy with long runs of
near-identical consecutive actions. A chunked policy conditioned on a short proprioceptive
history reaches the same fixed point there: it predicts the action it just executed. We
therefore require the demonstrations to be sparse in near-repeated actions,
\begin{equation}
  \Pr\bigl(\lVert a_{t+1} - a_t \rVert_\infty < \varepsilon_{\text{dw}}\bigr)
    \;\lesssim\; \tau_{\text{dw}} ,
  \label{eq:dwell}
\end{equation}
which the continuous approach profile~\eqref{eq:approach} satisfies by construction, and
which a waypoint-and-stop or minimum-jerk profile does not. The one place a deliberate
repeat is desirable is the end of the episode, where a small number of held frames is the
only supervision the policy receives for \emph{stopping} at the lift height; the trim
threshold in \S\ref{sec:method:demos} sets that count. Both~\eqref{eq:lead}
and~\eqref{eq:dwell} are checked automatically before training.

\subsection{Policy and training}
\label{sec:method:policy}

The policy is a denoising-diffusion behaviour-cloning model over action chunks. A PointNet
encoder maps the current cloud to a visual feature; it is concatenated with a short history
of the proprioceptive part of~\eqref{eq:obs} to form the conditioning vector. A residual MLP
denoiser predicts a chunk of $H$ future actions,
\begin{equation}
  \mathcal{L}(\phi)\;=\;
  \mathbb{E}_{(o,\,a_{t:t+H}),\;k,\;\epsilon}
  \Bigl[\;\bigl\lVert \epsilon - \epsilon_\phi\bigl(\tilde a^{(k)}_{t:t+H},\,k,\,c(o)\bigr)\bigr\rVert^2\Bigr],
  \label{eq:bc}
\end{equation}
with $k$ a sampled diffusion step, $\tilde a^{(k)}$ the noised chunk and $c(o)$ the
conditioning. An exponential moving average of the weights is retained for evaluation. All
checkpoints are kept and evaluated (\S\ref{sec:method:eval}) rather than selecting the last,
because the best checkpoint is consistently early in training.

\paragraph{Co-training with real demonstrations.}
Simulated demonstrations supply coverage; a small set of real teleoperated demonstrations
supplies the sensor and dynamics statistics of the target domain. The two are concatenated
without oversampling and normalised jointly, so a single policy sees both distributions with
their natural mixing ratio. Real demonstrations are recorded through the same shared
perception and action pipeline as the simulated ones, so the only differences between the two
data sources are those the hardware actually introduces.

\subsection{From specialist to generalist}
\label{sec:method:generalist}

Nothing in \S\ref{sec:method:synth}--\ref{sec:method:demos} is specific to one object
category: the planner requires a mesh and a material, and the executor requires a planned
grasp. A generalist is therefore obtained by extending the scene distribution rather than by
changing the method. Each category contributes a mesh set---scanned, or reconstructed from
photographs with an image-to-3D model, then rescaled to the category's physical extent---and a
material preset ($E$, $\nu$, $\rho$, yield); a scene draws a category, then a mesh within it,
then the usual scale, shape and material randomisation on top. Demonstrations from all
categories are pooled into a single dataset with per-category episode counts balanced by
construction, and one policy is trained on the pool.

Because the demonstrator is itself per-category, its statistics (planning success, induced
stress, contact area) are recorded alongside every episode. This matters for diagnosis: a
generalist that fails on a category may be failing because the policy lacks capacity across
categories, or because the demonstrator itself was weak on that geometry, and these two are
distinguishable only if the demonstrator is audited per category rather than in aggregate.

\subsection{Evaluation protocol}
\label{sec:method:eval}

Because the pipeline has many interacting choices, all simulated evaluation goes through one
harness with a fixed specification: a fixed number of episodes, a fixed number of parallel
environments, and a deterministic per-batch randomisation seed, so that environment $j$ of
batch $i$ presents an identical scenario to every checkpoint and every method. Geometry-level
randomisation is rebuilt on a fixed batch cadence rather than continuously, and any periodic
randomisation that a training server would otherwise apply is frozen for the duration of an
evaluation. We report task success, the fraction of episodes that ever reached the success
band (which separates grasping from holding failures), the gap between the two, and stress
percentiles over successful episodes. Deformable simulation on GPU is not bit-reproducible
across runs, so identical initial conditions still yield small variation in outcome; the
fixed scenario sequence controls the initial conditions, not the rollout noise.

On hardware, evaluation uses a fixed set of object placements inside the same workspace
rectangle used for training-time randomisation, with the object's orientation varied between
trials. Because a real trial cannot be repeated on an identical object instance---produce
varies, and a squeezed item is not reusable---real evaluation reports success over trials at
a stated sample size, and gentleness on hardware is assessed qualitatively unless a
destructive test is run.

\subsection{Real-robot deployment}
\label{sec:method:deploy}

Deployment runs the same policy, the same perception pipeline and the same action decoder as
training, with three additions.

\paragraph{Residual perception calibration.}
Camera extrinsics are calibrated once from fiducials; on top of that we apply a measured
residual translation. To obtain it we construct \emph{step-aligned real/simulated observation
pairs}: a recorded real episode's actions are replayed open-loop in simulation with the
object placed at the real first-frame fingertip position and the simulated home pose matched
to the real one, so that episode $i$, step $t$ is the same commanded state in both domains.
Because the proprioceptive channels of such pairs agree to within a few millimetres, any
residual displacement between the two point clouds is attributable to perception rather than
to control. Segmenting the cloud and measuring the displacement of the arm---whose pose is
pinned by proprioception---yields a systematic offset along the camera ray, which is then
applied as a constant correction to the camera translation. This procedure is a reusable
diagnostic: it converts ``the sim-to-real gap'' into a decomposition with a correctable part
(a rigid offset) and a residual part (sensor noise, geometry mismatch).

\paragraph{Execution.}
Action chunks are smoothed by an exponential moving average across chunk boundaries and
bounded by the same rate limit~\eqref{eq:rate} used during collection, so that an
out-of-distribution output is executed as a bounded walk. The arm is commanded in servo mode
at the control rate, and no tactile or force sensing is used at any point in the deployed
loop.

\subsection{Implementation details}
\label{sec:method:impl}

Table~\ref{tab:params} lists the parameters needed to reproduce the pipeline. Values are the
ones used for the results reported below; where a quantity is randomised, the range is given.

\begin{table}[t]
\centering\small
\begin{tabular}{ll}
\toprule
\textbf{Simulation} & MPM deformable, $\Delta t = 1/30$\,s, 235 substeps, grid density 250 \\
\quad material (nominal / DR) & $E$: 0.3 / $[0.2,0.3]$\,MPa; $\nu$: 0.35 / $[0.32,0.38]$;
                                $\rho$: $[900,1100]$\,kg\,m$^{-3}$; yield $\approx$ 40\,kPa \\
\quad coupling friction & $[3.5, 4.5]$ \\
\midrule
\textbf{Scene randomisation} & \\
\quad object position & $x\!\in\![0.29,0.48]$\,m, $y\!\in\![-0.11,0.11]$\,m (physical table rectangle) \\
\quad orientation & yaw $\mathcal{U}(-180^\circ,180^\circ)$; pitch/roll $\pm45^\circ$;
                    flip w.p.\ 0.25 to $\pm[160^\circ,180^\circ]$ \\
\quad scale / shape & scale $[0.8,1.5]$; bend $\pm25^\circ$; twist $\pm20^\circ$;
                      taper $\pm0.15$; one-axis scale $[0.95,1.15]$ \\
\quad arm home jitter & $\pm2$\,cm per axis \\
\midrule
\textbf{Perception} & crop $[0.20,-0.215,0.004]$--$[0.71,0.215,0.45]$\,m; FPS to $N_p\!=\!1024$ \\
\quad outlier filter (real only) & voxel 1\,cm, $\ge$23 neighbours \\
\quad arm-focus & full weight below $z\!=\!0.15$\,m or within $r\!=\!0.13$\,m of the EE;
                  else sampled at $0.15\times$ \\
\midrule
\textbf{Action} & position $[0.26,-0.225,0.003]$--$[0.59,0.225,0.50]$\,m; gripper $[0,0.088]$\,m \\
\quad rotation encoding & Euler \textsc{xyz} with frame offset $(180^\circ,0,0)$ \\
\quad rate limit $\rho$ & $(4.5, 4.5, 5.5)$\,mm, $(12, 12, 45)$\,mrad, $5$\,mm per step \\
\midrule
\textbf{Demonstrator} & CMA-ES, 1145 evaluations, multi-start; area floor 15\,mm$^2$
                        (scaled by scale$^2$) \\
\quad score weights & alignment, peak, pressure, azimuth bound $\bar\theta = 60^\circ$;
                      pose jitter $30^\circ$ \\
\quad execution phases & settle 1, close 20, firm 8, lift, hold; extra closure 5\,mm \\
\quad approach & $f\!\sim\!\mathcal{U}[0.45,0.75]$, $v_{\text{ref}}\!=\!2.4$\,mm/step,
                 $T_{\text{app}}\!\in\![40,130]$ steps \\
\quad post-processing & held runs $>$12 collapsed to 10 frames; pinch filter \\
\midrule
\textbf{Policy} & PointNet encoder $\rightarrow$ 512-d; MLP $[1024,1024,1024]$, residual \\
\quad chunking & horizon $H\!=\!4$, proprio history 2, one cloud per step, 20 denoising steps \\
\quad optimisation & batch 128, lr $10^{-4}$ cosine (100 warm-up), 600 epochs, EMA 0.995 \\
\quad size & $\approx$2.9\,M parameters \\
\midrule
\textbf{Control / hardware} & 7-DoF arm, parallel jaw, fixed external depth camera, 30\,Hz,
                              4 executed actions per chunk \\
\midrule
\textbf{Evaluation} & 200 episodes, 5 parallel environments, fixed base seed,
                      geometry rebuilt every 4 batches \\
\bottomrule
\end{tabular}
\caption{Pipeline parameters. Ranges denote per-scene randomisation.}
\label{tab:params}
\end{table}
